\documentclass[11pt]{amsart}
% !Mode:: "TeX:UTF-8"
\usepackage{amssymb,amsfonts,amsmath,amsthm}
\usepackage[all,arc]{xy}
\usepackage{enumerate}
\usepackage{mathrsfs}
\usepackage[toc,page]{appendix}
\usepackage[left=3cm, right=3cm, bottom=3cm]{geometry}
\usepackage{tabularx}
\usepackage{graphicx}
\usepackage{url}
\usepackage{color}
\usepackage{tikz-cd}  % graph
\usepackage{mathdots} %\ddots, \vdots, \iddots
\usepackage{romannum} %Roman number
\usepackage{hyperref} %hyperlink for reference
\newtheorem{thm}{Theorem}[section]
\newtheorem{qthm}[thm]{Quasi-Theorem}
\newtheorem*{thm*}{Theorem}
\newtheorem*{cor*}{Corollary}
\newtheorem*{prop*}{Proposition}
\newtheorem{cor}[thm]{Corollary}
\newtheorem{prop}[thm]{Proposition}
\newtheorem{lem}[thm]{Lemma}
\newtheorem{conj}[thm]{Conjecture}
\newtheorem{quest}[thm]{Question}

\theoremstyle{definition}
\newtheorem{defn}[thm]{Definition}
\newtheorem{defns}[thm]{Definitions}
\newtheorem{con}[thm]{Construction}
\newtheorem{exmp}[thm]{Example}
\newtheorem{exmps}[thm]{Examples}
\newtheorem{notn}[thm]{Notation}
\newtheorem*{notn*}{Notation}
\newtheorem{notns}[thm]{Notations}
\newtheorem{addm}[thm]{Addendum}
\newtheorem{exer}[thm]{Exercise}

\theoremstyle{remark}
\newtheorem{rem}[thm]{Remark}
\newtheorem{rems}[thm]{Remarks}
\newtheorem{warn}[thm]{Warning}
\newtheorem{sch}[thm]{Scholium}
\newtheorem*{idea*}{Idea}


\newcommand{\Spec}{{\rm Spec}}
%\newcommand{\ch}{{\rm ch}}

% Lei's command
\newcommand{\sshf}[1]{\mathcal{O}_{#1}} % structure sheaf
\newcommand{\shf}[1]{\mathscr{#1}} % sheaf
\newcommand{\proj}[1]{\textrm{proj}(#1)} % Proj
\newcommand{\prj}[1]{\mathbb{P}^{#1}} % projective space
\newcommand{\map}[3]{#1: #2\rightarrow#3}
\newcommand{\ts}[1]{\otimes_{#1}} % tensor
\newcommand{\iso}{\simeq} % isomorphism
\newcommand{\ses}[3]{0\rightarrow#1\rightarrow#2\rightarrow#3\rightarrow{0}} % short exact sequence
\newcommand{\ch}[1]{\textrm{ch}(#1)} %Character function%
\newcommand{\pf}[1]{\textrm{pf}(#1)}
\newcommand{\embedding}{\hookrightarrow}
\newcommand{\rk}[1]{\textrm{rank}(#1)} % rank
\newcommand{\is}[1]{\mathscr{I}_{#1}} % ideal sheaf
\newcommand{\lsp}[1]{\left\langle{#1}\right\rangle}
\newcommand{\dlc}[1]{\omega^{\bullet}_{#1}}
\newcommand{\spec}[1]{\textrm{spec}#1}

\makeatletter
\let\c@equation\c@thm
\makeatother
\numberwithin{thm}{section}
\numberwithin{equation}{section}

\bibliographystyle{plain}

\title[Student Research Program of mathematics]{Chevalley restriction theorem for classical simple Lie algebra}

\author{Jiahao Liang}

\begin{document}
\pagenumbering{arabic}
\maketitle
\begin{abstract}
In this article, I demonstrate the background knowledge 
,  which is mainly about Lie group-Lie algebra correspondence and Chevalley restriction theorem and prove that it holds for the case of classical simple Lie groups.
\end{abstract}
% keywords can be removed

\flushbottom

%\tableofcontents

%\newpage

\renewcommand{\thefootnote}{\fnsymbol{footnote}}
\footnotetext[1]{MSC2010 Class: 14D20, 14J60}
\footnotetext[2]{Key words:complex semisimple Lie algebra, Cartan subalgebra, Weyl group, invariant polynomial function}

\section{Introduction}%以下内容是介绍
At the International Congress of Mathematicians in 1950, Claude Chevalley proposed the famous Chevalley restriction theorem, which states that the conjugate invariant functions on a semisimple Lie algebra are isomorphic to the Weyl group invariant functions on a Cartan subalgebra. The theorem was fully proved by Robert Steinberg in 1965 and generalized to semisimple algebraic groups. Since then, mathematicians have made many advances in the higher-dimensional \cite{CN} generalization of Chevalley's restriction theorem. When we study the properties of high-dimensional Hitchin morphism, the Chevalley restriction theorem plays a key role.

In this article, I demonstrate the background knowledge and prove that it holds true for the case of classical simple Lie groups.
\subsection{Chevalley restriction theorem}
\begin{itemize}
    \item $G$: a \textbf{complex connected semi-simple} Lie group;
    \item $\mathfrak{g}$: the Lie algebra of $G$.
    \item $\mathbb{C} \left[ \mathfrak{g} \right] ^G$: the polynomial functions on $\mathfrak {g}$ which are invariant under the \textbf{adjoint $G$-action}.
    \item $\mathfrak{h}$: a \textbf{Cartan subalgebra} of $\mathfrak{g}$.
    \item $W$: the \textbf{Weyl group} of $G$.
    \item $\mathbb{C} \left[ \mathfrak{h} \right] ^W$: the polynomial functions on $\mathfrak{h}$ which are invariant under the natural action of $W$.
\end{itemize}
Then restriction of polynomial functions induces an isomorphism as \textbf{graded algebra}:
$$
\mathbb{C} \left[ \mathfrak{g} \right] ^G\cong \mathbb{C} \left[ \mathfrak{h} \right] ^W.
$$

\subsection{Elucidate} 
\begin{itemize}
    \item We introduce how to obtain the Lie algebra of a Lie group in the section of the Lie group–Lie algebra correspondence.
    \item Adjoint $G$-action (for linear Lie groups): We have a $G$-action on $\mathfrak{g}$ defined by $g.x = g^{-1}xg$. This induces an action on $\mathbb{C}[\mathfrak{g}]$, $g.P(x): = P(g^{-1}xg)$. In the section of the Lie group–Lie algebra correspondence, we will prove that this $G$-action is well - defined. $P$ is a polynomial function, which means that when we have chosen a basis, the function $P$ can be written as a polynomial of the coordinates. In summary, 
$$\mathbb{C} \left[ \mathfrak{g} \right] ^G:=\left\{ P:\mathfrak{g} \rightarrow \mathbb{C} \left| \begin{array}{c}
	\text{for all } (x,g)\in \mathfrak{g}\times G, P\left( g^{-1}xg \right) =P\left( x \right)\\
	P\left( \sum_{k=1}^n{x_ke_k} \right) =\sum{a_{k_1..k_n}x_{1}^{k_1}..x_{n}^{k_n}}\\
\end{array} \right. \right\} .
$$
\item For complex  semisimple Lie algebra $\mathfrak{g}$, we have root space decomposition \cite{GTM09}
$$\mathfrak{g} =\mathfrak{h} \oplus \bigoplus_{\alpha \in \Phi}^{}{\mathfrak{g} _{\alpha}},$$
while $\mathfrak{g} _{\alpha}:=\left\{ x\in \mathfrak{g} |\forall h\in \mathfrak{h},\mathrm{ad}\,h(x) =\alpha \left( h \right) x \right\} $ , $\alpha \in \mathfrak{h}^*$.  And the root system $\Phi$ has a base $\Delta$, which can span the whole space $\mathfrak{h}^*$.
\item The killing form $\kappa(x,y):=\mathrm{tr} (\mathrm{ad} x \, \mathrm{ad} y)$ and its restriction on $\mathfrak{h}$ is  non degenerate\cite{GTM09} , so for all 
$ \alpha \in \mathfrak{h}^*$, there exists unique $t_{\alpha}$ such that 
$$
\alpha =\kappa \left( t_{\alpha},- \right) :\mathfrak{h} \longrightarrow \mathbb{C},\quad h\longmapsto \kappa \left( t_{\alpha},h \right) .
$$
Let $E:=\mathbb{R} \alpha _1\oplus \mathbb{R} \alpha _1\oplus \cdots \oplus \mathbb{R} \alpha _l$, and $(\alpha,\beta):=\kappa(t_{\alpha},t_{\beta})$ then $\Phi$ becomes a root system of euclidean space\cite{GTM09} $E$. 
\item What is the Weyl group? Firstly, we define the Weyl group  \cite{GTM09} $W :=\left< \sigma _{\alpha}:\alpha \in \Phi \right>\subseteq O(E) $, while $\sigma_{\alpha} :E\to E $  such that  $$
\sigma _{\alpha}\left( \beta \right) =\beta -\frac{2\left( \beta ,\alpha \right)}{\left( \alpha ,\alpha \right)}\alpha .
$$
\item What is its natural action? Let $\Delta =\{\alpha _1,...,\alpha_l\}\subseteq \Phi$  be a base and fix a reflection $\sigma_{\alpha}:\alpha \in \Phi$, since $\sigma_{\alpha}(\Phi)=\Phi$, we have $
\Delta \hookrightarrow \Phi 
$. We denote $\varphi:\alpha \mapsto  t_{\alpha}$. Since $\Delta$ is a basis of $\mathfrak{h}^*$, there exists a unique $\tilde{\sigma}_{\alpha}$ that makes the following diagram commute. 
\begin{center}
\begin{tikzcd}
	\beta & \Delta & \Phi \\
	{t_{\beta}} & {\mathfrak{h}} & {\mathfrak{h}} 
	\arrow[maps to, from=1-1, to=2-1]
	\arrow["{\sigma_{\alpha}}", hook, from=1-2, to=1-3]
	\arrow["\varphi"', hook', from=1-2, to=2-2]
	\arrow["\varphi", hook', from=1-3, to=2-3]
	\arrow["{\tilde{\sigma}_{\alpha}}", dashed, from=2-2, to=2-3]
\end{tikzcd}
\end{center}
Thus, we have $W$ acts on $\mathfrak{h}$, extended by actions of its generators
$$
W\times \mathfrak{h} \longrightarrow \mathfrak{h} \quad \left( \sigma _{\alpha},\sum_i{c_it_{\alpha _i}} \right) \longmapsto \sum_i{c_i\tilde{\sigma}_{\alpha}\left( t_{\alpha _i} \right)}.
$$
In fact we can indentify $t_{\alpha}=\alpha $ via $\varphi$.
\end{itemize}

\vspace{0.5cm}
\section{Lie group-Lie algebra correspondence}
\subsection{Lie algebra of linear Lie group} 
 Let $\mathbb{K}=\mathbb{R} $ or $\mathbb{C}$. \begin{lem}[chap.3 in \cite{SGLG}]
 Let $x,y\in M_n(\Bbb{C})$ then we have \\
(i) (Trotter Product Formula)
$$
\underset{k\rightarrow \infty}{\lim}\left( \exp \frac{x}{k}\exp \frac{y}{k} \right) ^k=\exp \left( x+y \right) .
$$
(ii) (Commutator Formula) 
$$
\underset{k\rightarrow \infty}{\lim}\left( \exp \frac{x}{k}\exp \frac{y}{k}\exp \frac{-x}{k}\exp \frac{-y}{k} \right) ^{k^2}=\exp \left( xy-yx \right) .
$$
\end{lem}

\begin{defn}
 Let $
G\subseteq \mathrm{GL}_n\left( \mathbb{K}\right) $ be a closed subgroup of $\mathrm{GL}_n\left( \mathbb{K} \right) 
$, define $$
\mathbf{L}\left( G \right) :=\left\{ x\in M_n\left( \mathbb{K} \right) |\forall t\in \mathbb{R},\exp \left( t x \right) \in G \right\} 
$$
with Lie bracket $
\left[ x,y \right] :=xy-yx
$, then $\mathbf{L}(G)$  be come a Lie algebra, called the Lie algebra of Lie group.
\end{defn}

\begin{proof}
  Using Lemma 2.1 we know $\mathbf{L}(G)$ is a Lie algebra with $[x,y]=xy-yx$. 
(Refer to \cite{SGLG}Chap.3 for a detailed proof)    
\end{proof}


\begin{prop}[Chap.4 in \cite{SGLG}]
Let $\beta:V\times V\to \Bbb{K} $ be a bilinear map and $V$ be a linear space of finite dimension over $\Bbb{K}=\Bbb{R} $ or $\Bbb{C}$, then we have a closed subgroup 
$$
G=\mathrm{Aut}\left( V,\beta \right) :=\left\{ g\in \mathrm{GL}\left( V \right) :\forall v,w\in V,\beta \left( gv,gw \right) =\beta \left( v,w \right) \right\} 
$$
and its Lie algebra
$$
\\
\mathbf{L}G=\left\{ x\in \mathfrak{g} \mathfrak{l} \left( V \right) :\beta \left( xv,w \right) +\beta \left( v,xw \right) =0 \right\} .
$$
If $\beta \left( v,w \right) =v^\top Bw$ then
$$
\mathbf{L}G=\left\{ x\in \mathfrak{g} \mathfrak{l} \left( V \right) :x^\top B+Bx=0 \right\};
$$
If $\beta \left( v,w \right) =v^*Bw$ then
$$
\mathbf{L}G=\left\{ x\in \mathfrak{g} \mathfrak{l} \left( V \right) :x^*B+Bx=0 \right\}.
$$
\end{prop}

\begin{exmps}
Let $$
I_{p,q}:=\left( \begin{matrix}
	0&		I_p\\
	I_q&		0\\
\end{matrix} \right),\, J_n:=\left( \begin{matrix}
	0&		I_n\\
	-I_n&		0\\
\end{matrix} \right).
$$
By Proposition 2.4 
\begin{itemize}
    \item If $\mathrm{O}_{p,q}\left( \mathbb{K} \right) :=\left\{ x\in \mathrm{GL}_{p+q}\left( \mathbb{K} \right) :x^\top I_{p,q}x=I_{p,q} \right\} $ then
$$
\mathfrak{o} _{p,q}\left( \mathbb{K} \right) =\left\{ x\in \mathfrak{g} \mathfrak{l} _{p+q}\left( \mathbb{K} \right) |x^\top I_{p,q}+I_{p,q}x=0 \right\} =\mathbf{L}\mathrm{O}_{p,q}\left( \mathbb{K} \right) ;
$$
\item If$
\mathrm{sp}_n\left( \mathbb{K} \right) :=\left\{ g\in \mathrm{GL}_{2n}\left( \mathbb{K} \right) :g^\top J_n+J_ng=0 \right\} 
$ then 
\begin{align*}
    \mathfrak{s} \mathfrak{p} _n\left( \mathbb{K} \right) :&=\left\{ x\in \mathfrak{g} \mathfrak{l} _{2n}\left( \mathbb{K} \right) :x^\top J_n+J_nx=0 \right\} \\
 &=\left\{ \left( \begin{matrix}
	A&		B\\
	C&		D\\
\end{matrix} \right) \in \mathfrak{gl}_{2n}\left( \mathbb{K} \right) :B^\top =B,C^\top =C,D=-A^\top \right\} ;
\end{align*}
\item Let $G:=\mathrm{SO}_{p,q}(\Bbb{C}):=\mathrm{O}_{p,q}(\Bbb{C})\cap \mathrm{SL}_n(\Bbb{C})$ , then we have
$$
\mathfrak{so} _{p,q}(\Bbb{C}):=\mathbf{L}\mathrm{SO}_{p,q}\left( \mathbb{C} \right) =\left\{ x\in \mathfrak{g} \mathfrak{l} _{p+q}\left( \mathbb{C} \right) |x^\top I_{p,q}+I_{p,q}x=0 \right\} =\mathfrak{o} _{p,q}.
$$
\end{itemize}
\end{exmps}

\subsection{The Lie functor}
\begin{prop}
 Let $\varphi: G_1\to G_2$ be a morphism of Lie groups, then there exists unique morphism of Lie algebras $\mathbf{L}(\varphi)$ making the following diagram commute
\begin{center}   
\begin{tikzcd}   
G_1 \arrow[dd,leftarrow, no head,"\mathrm{exp}_{G_1}"'] \arrow[r,"\varphi"] & G_2 \arrow[dd,leftarrow, no head,"\mathrm{exp}_{G_2}"]\\   
&\\   
\mathbf{L}(G_1)\arrow[r,dashed,"\mathbf{L}(\varphi)"] & \mathbf{L}(G_2)   
\end{tikzcd} 
\end{center}
And we have a functor $\mathbf{L}$ from the category of linear Lie groups to the category of Lie algebras over $\mathbb{K}$.
\end{prop}

\begin{proof}
(Refer to \cite{SGLG} Chap.3 for a detailed proof)   For the existence of $\mathbf{L}(\varphi)$, we define
$$
\mathbf{L}\left( \varphi \right) :\mathbf{L}(G_1)\longrightarrow \mathbf{L}(G_2)\qquad x\longmapsto \left. \frac{\mathrm{d}}{\mathrm{d}t} \right|_{t=0}\varphi \left( \exp _{G_1}\left( tx \right) \right) .
$$
\end{proof}

\begin{exmps}[\textbf{Adjoint G-action}]
For linear Lie groups, let $g\in G$, then we have an isomorphism of Lie groups:
$$
c_g:G\longrightarrow G,\,\, h\longmapsto ghg^{-1}.
$$
Then it is easy to know 
$$
L\left( c_g \right) :\mathbf{L}(G)\longrightarrow \mathbf{L}(G), \,\,x\longmapsto gxg^{-1}.
$$
which is an isomorphism of Lie algebras. This induces a smooth action of the Lie group on $\mathbf{L}(G)$
$$
G\times \mathbf{L}G\longrightarrow \mathbf{L}G\quad \left( g,x \right) \longmapsto gxg^{-1}.
$$
\end{exmps}

\subsection{Adjoint Representation}. 
\begin{lem}
Let $x,y\in \mathfrak{gl}(V)$ be two arbitrary linear maps, then we have
$$
e^xye^{-x}=e^{\mathrm{ad}\,x}y.
$$
\end{lem}
\begin{proof}
Let $\lambda _x:y\longmapsto xy,\,\rho _x=y\longmapsto yx
$, which are commutative, then we have
$$
e^xye^{-x}=e^{-\rho _x}e^{\lambda _x}y=e^{\lambda _x-\rho _x}y.
$$
\end{proof} 

\begin{prop}
We define the adjoint representation of $G$ by
$$
\mathrm{Ad}:G\longrightarrow \mathrm{Aut}\left( G \right) ,\quad g\longmapsto c_g.
$$
Then we have
$$
\mathbf{L}\left( \mathrm{Ad} \right) =\mathrm{ad}:\mathbf{L}G\longrightarrow \mathbf{L}\mathrm{Aut}\left( G \right) \subseteq \mathfrak{g} \mathfrak{l} \left( \mathbf{L}G \right).$$
\end{prop}

\begin{proof}
 Since $e^xye^{-x}=e^{\mathrm{ad}\,x}y$, we have the following commutative diagram
\begin{center}
\begin{tikzcd}
	G && {\mathrm{Aut} (G)} \\
	\\
	{\mathbf{L}G} && {\mathfrak{gl}(\mathbf{L}G)}
	\arrow["{\mathrm{Ad}}", from=1-1, to=1-3]
	\arrow["{\exp_1}", from=3-1, to=1-1]
	\arrow["{\exp_2}"', from=3-3, to=1-3]
	\arrow["{\mathrm{ad}}"', from=3-1, to=3-3]
\end{tikzcd}
\end{center}
For uniqueness we know $\mathbf{L}\left( \mathrm{Ad} \right) =\mathrm{ad}$. 
\end{proof}
 

\begin{defn}[Inner isomorhism of $\mathfrak{g}$]
For the linear Lie group over $\Bbb{C}$, define
$$
\mathrm{Inn}\left( \mathfrak{g} \right) :=\left< e^{\mathrm{ad}\,x}:x\in \mathfrak{g} \right> =\left< \mathrm{Ad}\left( e^{-x} \right) :x\in \mathfrak{g} \right> .
$$
\end{defn}

\begin{rem}
The group $\mathrm{Inn}\left( \mathfrak{g} \right)$ acting on $\mathfrak{g}$ can identify as some adjoint $G$-actions because of
$$
\mathrm{Inn}\left( \mathfrak{g} \right) \subseteq \left\{ \mathrm{Ad}\left( g \right) :g\in G \right\}.
$$
\end{rem}

\begin{cor*}
 If $H$ is a normal Lie subgroup of $G$, then $\mathbf{L}H$ is an ideal of $\mathbf{L}G$.
\end{cor*}
\begin{proof}
     $H\lhd G$ implies that $\mathrm{Ad}:G\to \mathrm{Aut}(H)$, so we have $\mathrm{ad}:\mathbf{L}G \to \mathfrak{gl}(\mathbf{L}H)$. So for all $x\in \mathbf{L}G,\mathrm{ad}\,x$ preserves $\mathbf{L}H$.
\end{proof}

\vspace{0.5cm}
\section{Classification of complex Simple Lie algebra and their Weyl groups}

The materials in the section are well-known and readers should refer to \cite{GTM09} for more details. 

\begin{lem}\label{lem:1}
We have a canonical isomorphism as Lie groups
$$
\mathrm{SO}_{n,n}\left( \mathbb{C} \right) \cong \mathrm{SO}_{2n}\left( \mathbb{C} \right) .
$$
In fact $\mathrm{SO}_{n,n}\left( \mathbb{C} \right) \to \mathrm{SO}_{2n}\left( \mathbb{C} \right) ,\,\, g \mapsto ugu^{-1}$ gives an isomorphism with 
$$
u=\left( \begin{matrix}
	I_n&		I_n\\
	\mathrm{
}-iI_n&		iI_n\\
\end{matrix} \right) =\left( \begin{matrix}
	I&		0\\
	0&		iI\\
\end{matrix} \right) \left( \begin{matrix}
	I&		I\\
	-I&		I\\
\end{matrix} \right) ,\quad i^2=-1.
$$
\end{lem}
\textbf{Proof}: Calculate and verify directly.

\begin{rem}
Using the Lie functor, we have $\mathfrak{so}_{n,n}\left( \mathbb{C} \right) \cong \mathfrak{so}_{2n}\left( \mathbb{C} \right) $ as Lie algebras.
\end{rem}


\begin{lem}[The classification theorem \cite{GTM09} for semisimple Lie algebras]
      Every finite-dimensional semisimple Lie algebra over $\Bbb{C}$  is isomorphic to a direct sum of simple Lie algebras. These simple Lie algebras can be classified into the following types: $A_n$, $B_n$, $C_n$, $D_n$, $E_6$, $E_7$, $E_8$, $F_4$, and $G_2$. 
The general forms of these Lie algebras are given by: 
\begin{itemize}
    \item $A_n$ ($n \geq 1$): the Lie algebra $\mathfrak{sl}_{n+1}( \mathbb{C})$, with complex  Lie group $\mathrm{SL}_{n+1}(\Bbb{C})$ ;
    \item $B_n$ ($n \geq 2$): the Lie algebra $$\mathfrak{so}_{2n}(\mathbb{C})\simeq\mathfrak{so}(\mathbb{C}^{2n+1}, T) := \{x \in \mathfrak{gl}_{2n+1}(\mathbb{C}) : x^\top T + T x = 0\};$$ while $T := 
\begin{pmatrix}
 I_{n,n} & 0 \\
0 & 1 \\
\end{pmatrix}
\in M_{2n+1}(\mathbb{C})$
    \item $C_n$ ($n \geq 3$): 
the Lie algebra $$\mathfrak{g=sp}_{2n}(\mathbb{C})= \{x \in \mathfrak{gl}_{2n}(\mathbb{C}) : x^\top J_n + J_nx = 0 \};$$
    \item $D_n$ ($n \geq 4$):
 the Lie algebra $$\mathfrak{so}_{2n}(\mathbb{C})\simeq\mathfrak{so}(\mathbb{C}^{2n}, I_{n,n}):= \{x \in \mathfrak{gl}_{2n}(\mathbb{C}) : x^\top I_{n,n} + I_{n,n} x = 0\};$$
    \item $E_6$: a complex Lie algebra of dimension 78;
    \item $E_7$: a complex Lie algebra of dimension 133;
    \item $E_8$: a complex Lie algebra of dimension 248;
    \item $F_4$: a complex Lie algebra of dimension 52;
    \item $G_2$: a complex Lie algebra of dimension 14.
\end{itemize}
\end{lem}

\begin{defn}[Cartan subalgebra]
We say that the subalgebra $\mathfrak{h}$ is a Cartan subalgebra of $\mathfrak{g}$ if $\mathfrak{h}$ is nilpotent and $N_{\mathfrak{g}}(\mathfrak{h})=\mathfrak{h}$.
\end{defn}

\begin{lem}[Chap.6 in \cite{SGLG} ]
Let $\mathfrak{g}$ be a simisimple Lie algebra over $\Bbb{C}$, then the following are equivalent:
\begin{itemize}
\item $\mathfrak{h}$ is a Cartan subalgebra.
\item $\mathfrak{h}$ is a maximal torus subalgebra.
\item $\mathfrak{h}$ is a maximal abelian subalgebra consisting of semisimple elements.
\end{itemize}
\end{lem}


\begin{lem}[\cite{SGLG}Cartan subalgebra of classical simple Lie algebra]
Let 
\begin{enumerate}
\item $\mathfrak{h}$ denote one of Cartan subalgebra of $\mathfrak{g}$;
\item  $\varepsilon_j : \mathfrak{h} \to \mathbb{C}$ be the linear functions defined by 
$$\varepsilon_k(\mathrm{diag}(h, -h)) := h_k,\quad h=(h_1,...,h_n);$$
\item  $\Delta$ be the basis of the root system; 
\item $E_{i,j}$ denote the elementry matrix and $E_{k}:=E_{kk}$ .
\end{enumerate}
 Then we have we have the following results.
\begin{itemize}
\item $A_n$ ($n\geq1$):When $\mathfrak{g=sl}_{n + 1}(\mathbb{C})$ 
$$\mathfrak{h}=\mathrm{span}\{E_{j} - E_{j+1} : j = 1, \ldots, n\} = \{\mathrm{diag}(h, -\mathrm{tr}\,h) : h\in \mathbb{C}^n\}.$$
The corresponding root system and one of its bases are
$$A_{n-1}:= \{\varepsilon_j - \varepsilon_k :1\le j, k \ne  1\neq n-1 \};
$$
$$
\Delta =\{\varepsilon _j-\varepsilon _{j+1}|j=1,...,n\}.
$$
\item $B_n$ ($n\geq2$): When $\mathfrak{g=so}_{2n + 1}(\mathbb{C})$
$$\mathfrak{h} = \mathrm{span}\{E_{j} - E_{n+j} : j = 1, \ldots, n\} = \{\mathrm{diag}(h, -h, 0) : h \in \mathbb{C}^n\}.$$
The corresponding root system and one of its bases are
\[
B_n := \{\pm \varepsilon_j, \pm \varepsilon_j \pm \varepsilon_k : j, k = 1, \ldots, n, j \neq k\};
\]
$$
\Delta =\{\varepsilon _j-\varepsilon _{j+1}|j=1,...,n-1\}\cup \left\{ \varepsilon _n \right\} .
$$
\item $C_n$ ($n\geq3$): When $\mathfrak{g=sp}_{2n}(\mathbb{C})$
$$\mathfrak{h} = \mathrm{span}\{E_{j} - E_{n+j} : j = 1, \ldots, n\} = \{\mathrm{diag}(h, -h) : h \in \mathbb{C}^n\}.$$
The corresponding root system and one of its bases are
$$C_n := \{\pm 2\varepsilon_j, \pm \varepsilon_j \pm \varepsilon_k : j, k = 1, \ldots, n, j \neq k\};$$
$$
\Delta =\{\varepsilon _j-\varepsilon _{j+1}|j=1,...,n-1\}\cup \left\{ 2\varepsilon _n \right\} .
$$
\item $D_n$ ($n\geq4$): When $\mathfrak{g=so}_{2n}(\mathbb{C})$
$$\mathfrak{h} = \mathrm{span}\{E_{j} - E_{n+j} : j = 1, \ldots, n\} = \{\mathrm{diag}(h, -h) : h \in \mathbb{C}^n\}.$$
The corresponding root system and one of its bases are
$$D_n := \{\pm \varepsilon_j \pm \varepsilon_k : j, k = 1, \ldots, n, j \neq k\};$$
$$
\Delta =\{\varepsilon _j-\varepsilon _{j+1}|j=1,...,n-1\}\cup \left\{ \varepsilon _{n-1}+\varepsilon _n \right\} .
$$
\end{itemize}
\end{lem}

\begin{proof}
 The following proof is based on the example of type $A_n$ case and $C_n$ case,  as the proofs for other types are similar. 

The $A_{n-1}$ case  : It is easy to know $\mathfrak{h}=\{\mathrm{diag}(a_1,...,a_n)|a_1+...+a_n=0\}$ and
$$
\left( \mathrm{ad}h \right)\left( E_{ij} \right) =\left( h_i-h_j \right) E_{ij},\quad h=\mathrm{diag}(h_1,...,h_n).
$$
Further more
$$\mathfrak{s} \mathfrak{l} _n=\mathfrak{h} \oplus \bigoplus_{i\ne j}{\mathbb{C} E_{ij}}.$$
Then we can work out the corresponding root system is
$$A_{n-1}:= \{\varepsilon_j - \varepsilon_k :1\le j, k \ne  1\neq n-1 \};$$

For $C_n$ case, using
$$\begin{pmatrix}
a & b \\
c & d
\end{pmatrix}
\in \mathfrak{sp}_n(\mathbb{C}) \quad \iff \quad
d = -a^\top, \, b^\top = b, \, c^\top = c,$$
we see that $\mathfrak{g} := \mathfrak{sp}_n(\mathbb{C})$ has an abelian subalgebra
$$\mathfrak{h} = \mathrm{span}\{E_{j} - E_{n+j} : j = 1, \ldots, n\} = \{\mathrm{diag}(h, -h) : h \in \mathbb{C}^n\}.$$
If $\begin{pmatrix}
a & b \\
c & d
\end{pmatrix}
\in Z_{\mathfrak{g}}(\mathfrak{h})$, then for all $h=\mathrm{diag}(h_1,...,h_n)\in M_n(\Bbb{C})$
$$
0=\left( \begin{matrix}
	a&		b\\
	c&		d\\
\end{matrix} \right) 
\left( \begin{matrix}
	h&		\\
	&		-h\\
\end{matrix} \right) -
\left( \begin{matrix}
	h&		\\
	&		-h\\
\end{matrix} \right) 
\left( \begin{matrix}
	a&		b\\
	c&		d\\
\end{matrix} \right) 
=
\left( \begin{matrix}
	ah-ha&		-bh-hb\\
	ch+hc&		hd-dh\\
\end{matrix} \right) .
$$
$ah=ha$ implies that $a$ is a diagonal matrix. $hb+bh=0$ implies that $b=0$. So $Z_{\mathfrak{g}}(\mathfrak{h})=\mathfrak{h}$. 
Similarly, some of $E_{j,n+k}+E_{k,n+j}$ and $E_{n+k,j}+E_{n+j,k}$ form a basis of all $\mathfrak{g}_{\alpha}$. 

Then we can work out the corresponding root system is
$$C_n := \{\pm 2\varepsilon_j, \pm \varepsilon_j \pm \varepsilon_k : j, k = 1, \ldots, n, j \neq k\}.$$
\end{proof}


\begin{lem}[Generators of Weyl group]

Let $\Delta$ be a base of $\Phi$, then (refer to \cite{GTM09}) the Weyl group 
$$W :=\left< \sigma _{\alpha}:\alpha \in \Phi \right> =\left< \sigma _{\alpha}:\alpha \in \Delta \right>$$
\end{lem}

\begin{lem}\label{lem:2}
If the following conditions hold: 
\begin{enumerate}
    \item $E_1$ is a subspace of $E_2$ as a Euclidean space such that $E_2 = E_1 \oplus E_1^{\bot}$, when $E_1^{\bot} = \mathbb{R} \varepsilon$.
    \item  $H_1 $ subgroup of $O(E_1)$ and $H_2$ subgroup of $O(E_2)$.
    \item  For all $\sigma \in H_2, \sigma(\varepsilon) = \varepsilon$.
    \item $H_1 = \{\sigma|_{E_1} : \sigma \in H_2\}$.
\end{enumerate}
Then $H_1 \simeq H_2$.
\end{lem}
\begin{proof}
There is a surjective homomorphism $\varphi : H_2 \longrightarrow H_1,\,\, \sigma \longmapsto \sigma|_{E_1}$.
Its kernel is $$\mathrm{Ker}(\varphi) = \{\sigma \in H_2 : \sigma|_{E_1} = \mathrm{id}_{E_1}, \sigma(\varepsilon) = \varepsilon\} = \{\mathrm{id}\}.$$
Thus, it is an isomorphism.
\end{proof}

\begin{lem}[\cite{LA}Weyl group of classical simple Lie algebra]
Let $W$ be the Weyl group of $\mathfrak{g}$, then we have the following results.
\begin{itemize}
\item Type $A_{n-1}$ case: $W\simeq S_n$ with natural action ($E_j:=E_{jj}$) generated by
$$\sigma _{\varepsilon_k-\varepsilon_{k+1}}\left( E_{j} \right) =E_{\pi \left( j \right)},\pi =\left( k,k+1 \right) \in S_n.$$
\item Type $B_n$ case: $W\simeq S_n \ltimes \mathbb{Z}^{n}$ with natural action ($e_j:=E_{jj}-E_{n+j,n+j}$) generated by 
$$\sigma _{\varepsilon_k-\varepsilon_{k+1}}\left( e_{j} \right) =e_{\pi \left( j \right)},\pi =\left( k,k+1 \right) \in S_n;$$
$$
\sigma _{\varepsilon _n}\left( e_n\right) =-e_n.
$$
\item Type $C_n$ case:  $W\simeq S_n \ltimes \mathbb{Z}^{n}$ with natural action ($e_j:=E_{jj}-E_{n+j,n+j}$) generated by
$$\sigma _{\varepsilon_k-\varepsilon_{k+1}}\left( e_{j} \right) =e_{\pi \left( j \right)},\pi =\left( k,k+1 \right) \in S_n;$$
$$
\sigma _{2\varepsilon _n}\left( e_n \right) =-e_n.
$$
\item Type $D_n$ case:  $W\simeq S_n \ltimes \mathbb{Z}^{n-1}$ with natural action ($e_j:=E_{jj}-E_{n+j,n+j}$) generated by
$$\sigma _{\varepsilon_k-\varepsilon_{k+1}}\left( e_{j} \right) =e_{\pi \left( j \right)},\pi =\left( k,k+1 \right) \in S_n;$$
$$
\sigma _{\varepsilon _n+\varepsilon _{n-1}}\left( e_j\right) =-e_j, \quad j\in\{n-1,n\}.
$$
\end{itemize}
\end{lem}

\begin{proof}
 Type $A_{n-1}$ case: We claim that $$\kappa \left( x,y \right) =2n\mathrm{tr}\left( xy \right) ,\forall x,y\in \mathfrak{h}$$
Let $h=\mathrm{diag}(h_1,...,h_n)\in M_n(\Bbb{C})$, for
$\left( \mathrm{ad}h \right) ^2\left( E_{ij} \right) =\left( h_i-h_j \right) ^2E_{ij}$
we know 
$$
\begin{aligned}
	\kappa \left( h,h \right) &=\sum_{i\ne j}{\left( h_i-h_j \right) ^2}=2\sum_{i<j}{\left( h_i-h_j \right) ^2}\\
	&=2\sum_{i<j}{\left( h_i-h_j \right) h_i}+2\sum_{i<j}{\left( h_j-h_i \right) h_j}\\
	&=2\sum_{i=1}^{n-1}{h_i\sum_{j=i+1}^n{\left( h_i-h_j \right)}}+2\sum_{j=2}^n{h_j\sum_{i=1}^{j-1}{\left( h_j-h_i \right)}}\\
	&=2\sum_{i=1}^n{h_i\sum_{j=i}^n{\left( h_i-h_j \right)}}+2\sum_{i=1}^n{h_i\sum_{j=1}^{i-1}{2\left( h_i-h_j \right)}}\\
	&=\sum_{i=1}^n{h_i\sum_{j=1}^n{2\left( h_i-h_j \right)}}=2n\sum_{i=1}^n{h_{i}^{2}}=2n\mathrm{tr}\left( h^2 \right).
\end{aligned}
$$
By  $\kappa \left( x,y \right) =\frac{1}{4}\left[ \kappa \left( x+y,x+y \right) -\kappa \left( x-y,x-y \right) \right] $, our claim holds. 
More precisely 
$$
\sigma _{\alpha}\left( E_j \right) =e_j-\frac{2\mathrm{tr}\left( E_jt_{\alpha} \right)}{\mathrm{tr}\left( {t_{\alpha}}^2 \right)}t_{\alpha},\quad t_{\varepsilon _i}=E_i-E_{n+1}.
$$
Then we can work out the action, and conclude that
$$H_2:=\left< \sigma _{e_k-e_{k+1}}|k=1,...,n-1 \right> 
\simeq \left< \left( k,k+1 \right) \in S_n|k=1,...,n-1 \right> =S_n, $$
and let $\varepsilon  =e_1+e_2+....+e_n,E_1=\varepsilon ^{\bot}=\mathfrak{h}$, by \ref{lem:2} we have
$$W =\left\{ \sigma |_{E_1}:\sigma \in H_2 \right\}\simeq H_2\simeq S_n.$$
The type $D_n,B_n$ case is similar to $C_n$ case.  Let us consider the $C_n$ case and identify $\mathfrak{g=sp}_n(\mathbb{C}) $ as a subalgebra of $\mathfrak{sl}_{2n}(\mathbb{C})$, the killing form $$\kappa \left( x,y \right) =4n\mathrm{tr}\left( xy \right) ,\,\forall x,y\in \mathfrak{h}$$
then we can work out all reflections in a similar way.
\end{proof}

\vspace{0.5cm}
\section{Chevalley restriction theorem}
In the section of Introduction, we have stated the Chevalley restriction theorem.
\begin{thm}[Chevalley restriction theorem \cite{RTCG} ]
The restriction of polynomial functions induces an isomorphism as \textbf{graded algebra}:
$$
\mathbb{C} \left[ \mathfrak{g} \right] ^G\cong \mathbb{C} \left[ \mathfrak{h} \right] ^W
\qquad P\longmapsto P | _{\mathfrak{h}}.
$$
\end{thm}

\begin{prop}[Type $A_{n-1}$ case]
When $G=\mathrm{SL}_n,\,\mathfrak{g=sl}_n$, the induced map is an isomorphism.
\end{prop}
\begin{proof}
We have known that $\mathfrak{h}=\{\mathrm{diag}(a_1,...,a_n)|a_i\in\Bbb{C} ,a_1+...+a_n=0\}$. Let $E_{i,j}$ denote the elementry matrix and $E_{k}:=E_{kk}$. Then polynomial function means  
$$P\left( \displaystyle\sum_{i\ne j}{x_{i,j}E_{ij}+\displaystyle\sum_{k=1}^n{x_k}E_{kk}} \right) $$  are polynomials of $x_{i,j} $ and $x_k$ with the equation $$x_1+x_2+....+x_n=0.$$  Then we can identify the restriction map as 
$$
P\left( \sum_{i\ne j}{x_{i,j}E_{ij}+\sum_{k=1}^n{x_k}E_{kk}} \right) \longmapsto P|_{\mathfrak{h}}=P\left( \sum_{i\ne j}{0E_{ij}+\sum_{k=1}^n{x_k}E_{kk}} \right).
$$
Obviously the induced map is a homomorphism of graded algebra, the proof of this proposition is mainly divided into the following steps

\textbf{Step 1:} $\mathbb{C} \left[ \mathfrak{h} \right] ^W=\mathbb{C} \left[ \sigma _2,...,\sigma _n \right] $, while $\sigma_1,\sigma_2,...,\sigma_n$  are  elementary symmetric polynomials with equation $\sigma_1=0.$
\begin{proof}
 The action $W=S_n$ acting on $\mathbb{C} \left[ \mathfrak{h} \right]
$, is defined by
$$\pi .P\left( \displaystyle\sum_{k=1}^n{x_k}E_{kk} \right) =P\left( \displaystyle\sum_{k=1}^n{x_{\pi \left( k \right)}}E_{kk} \right) , \quad \pi \in S_n. $$
According to basic theorem of elementary symmetric polynomial, we have $$\mathbb{C} \left[ \mathfrak{h} \right] ^W\subseteq \mathbb{C} \left[ \sigma _1,\sigma _2,...,\sigma _n \right].$$
For $\sigma_1=0$,  $\mathbb{C} \left[ \mathfrak{h} \right] ^W=\mathbb{C} \left[\sigma _2,...,\sigma _n \right].$
\end{proof}

\textbf{Step 2:}  $\mathbb{C} \left[ \tau _2,...,\tau _n \right] \subseteq \mathbb{C} \left[ \mathfrak{g} \right] ^G$, while  $\tau_2,..,\tau_n$ are some  polynomials such that
$$
\det \left( tI_n-\sum_{i\ne j}{x_{i,j}E_{ij}-\sum_{k=1}^n{x_k}E_{kk}} \right) =t^n-0t^{n-1}+\tau _2t^{n-2}-\tau _3t^{n-3}+...+\left( -1 \right) ^n\tau _n.
$$
\begin{proof}
 It is welled known that $$\det \left( tI-g^{-1}Ag \right) =\det \left( tg^{-1}(I-A)g \right) =\det \left( tI-A \right). $$
Thus, the polynomial coefficients are invariant under the adjoint $G$-action,  which implies that they are invariant polynomials. 
\end{proof}

\textbf{Step 3:} The induced map to $\Bbb{C}\left[ \mathfrak{h} \right] ^W$ is well defined.  
\begin{proof}
We can identify the restriction map as
$$
P\left( \sum_{i\ne j}{x_{i,j}E_{ij}+\sum_{k=1}^n{x_k}E_{kk}} \right) \longmapsto P|_{\mathfrak{h}}=P\left( \sum_{i\ne j}{0E_{ij}+\sum_{k=1}^n{x_k}E_{kk}} \right).
$$
So if $P\in \mathbb{C} \left[ \mathfrak{g} \right] ^G$, then
$$
P\left( g^{-1}Ag \right) =P\left( A \right) \Rightarrow P\left( \sum_{i\ne j}{x_{i,j}g^{-1}E_{ij}g+\sum_{k=1}^n{x_k}g^{-1}E_{kk}g} \right) =P\left( \sum_{i\ne j}{x_{i,j}E_{ij}+\sum_{k=1}^n{x_k}E_{kk}} \right).
$$
Let $x_{i,j}=0$,  the equation still holds, so we have $P|_{\mathfrak{h}} \in \mathbb{C} \left[ \mathfrak{h} \right] ^G$.
According to the well-known linear algebra conclusion, these two matrices are conjugate 
$$
\sum_{k=1}^n{x_{\pi \left( k \right)}}E_{kk}\,\,,\sum_{k=1}^n{x_k}E_{kk}.
$$
So we have $P|_{\mathfrak{h}}\in \mathbb{C} \left[ \mathfrak{h} \right] ^G\subseteq \mathbb{C} \left[ \mathfrak{h} \right] ^W$. 
\end{proof}

\textbf{Step 4(Surjective)}:
\begin{proof}
In the equation 
$$
\det \left( tI_n-\sum_{i\ne j}{x_{i,j}E_{ij}-\sum_{k=1}^n{x_k}E_{kk}} \right) =t^n-0t^{n-1}+\tau _2t^{n-2}-\tau _3t^{n-3}+...+\left( -1 \right) ^n\tau _n.
$$
Let $x_{i,j}=0$, we have $\tau _i|_{\mathfrak{h}}=\sigma _i \forall i=2,...,n$ , which generate $\Bbb{C}[\mathfrak{h}]^W$. Obviously
$$
P\left( \tau _2,...,\tau _n \right) |_{\mathfrak{h}}=P\left( \tau _2|_{\mathfrak{h}},...,\tau _n|_{\mathfrak{h}} \right) =P\left( \sigma _2,...,\sigma _n \right).
$$
Thus $\mathrm{Im}\, \varphi=\Bbb{C}[\mathfrak{h}]^W$.  
\end{proof}

\textbf{Step 5(Injective)}: It suffices to show that $\mathrm{Ker}\,\varphi=0$ . 

\begin{proof}
Let $P\in\mathrm{ker}\,\varphi$, then $P(\mathfrak{h})=\{0\}$ and  for all $(x,g)\in\mathfrak{g}\times G$ 
$$P(g^{-1}xg)=P(x).$$
We denote 
$$
\mathfrak{g}^s=\{g^{-1}hg|g\in \mathrm{GL}_n,\,h\in\mathfrak{h}\}=\{g^{-1}hg|g\in \mathrm{SL}_n,\,h\in\mathfrak{h}\}
$$
the set of diagonalizable elements. Then we have $P(\mathfrak{g}^s)=\{0\}$. For $P$ is a continuous function, it suffices to show that $\mathfrak{g}^s$ is dense in $\mathfrak{g}$ .

Let $x \in \mathfrak{g}$ , we construct  a sequence $\{x_k\}$ in $\mathfrak{g}^s$ such that $\displaystyle\lim_{k\to \infty}x_k= x.$ We know there exists $g\in G$ such that 
$$
g^{-1}xg=\left( \begin{matrix}
	J\left( \lambda _1,n_1 \right)&		&		&		\\
	&		J\left( \lambda _2,n_2 \right)&		&		\\
	&		&		\ddots&		\\
	&		&		&		J\left( \lambda _r,n_r \right)\\
\end{matrix} \right),\quad J\left( \lambda_j,n_j \right) =\left( \begin{matrix}
	\lambda_j&		1&		&		&		\\
	&		\lambda_j&		1&		&		\\
	&		&		\ddots&		&		\\
	&		&		&		\lambda_j&		1\\
	&		&		&		&		\lambda_j\\
\end{matrix} \right) _{n_j\times n_j}.
$$
So let
$$
g^{-1}x_kg=\left( \begin{matrix}
	A_{k,1}&		&		&		\\
	&		A_{k,2}&		&		\\
	&		&		\ddots&		\\
	&		&		&		A_{k,r}\\
\end{matrix} \right),\quad A_{k,j}=\left( \begin{matrix}
	\lambda_j&		1&		&		&		\\
	&		\lambda_j&		1&		&		\\
	&		&		\ddots&		&		\\
	&		&		&		\lambda_j&		1\\
	\frac{1}{k}&		&		&		&		\lambda_j\\
\end{matrix} \right) _{n_j\times n_j},
$$
then for all $j$, we have $$\underset{k\rightarrow \infty}{\lim}A_{k,j}=J\left( \lambda_j ,n_j \right).$$ Thus 
$$
\underset{k\rightarrow \infty}{\lim}x_k=\underset{k\rightarrow \infty}{\lim}g\left( \begin{matrix}
	A_{k,1}&		&		\\
	&		\ddots&		\\
	&		&		A_{k,r}\\
\end{matrix} \right) g^{-1}=g\left( \begin{matrix}
	J\left( \lambda _1,n_1 \right)&		&		\\
	&		\ddots&		\\
	&		&		J\left( \lambda _r,n_r \right)\\
\end{matrix} \right) g^{-1}=x,
$$
and the characteristic polynomials 
$$
\det \left( \lambda I-A_{k,j} \right) =\left| \begin{matrix}
	\lambda -\lambda _j&		-1&		&		&		\\
	&		\lambda -\lambda _j&		-1&		&		\\
	&		&		\lambda -\lambda _j&		&		\\
	&		&		&		\ddots&		-1\\
	\frac{1}{k}&		&		&		&		\lambda -\lambda _j\\
\end{matrix} \right|_{n_j\times n_j}=\left( \lambda -\lambda _j \right) ^{n_j}-\frac{1}{k}
$$
do not have multiple roots, so $A_{k,j}$ is diagonalizable, so does $x_k$. This proof ends.   
\end{proof}

In summary, we have proved that  the induced map is  an isomorphism. 
\end{proof}

\begin{cor*}[Generators]
When $G=\mathrm{SL}_n,\,\mathfrak{g=sl}_n$, we have found out all generators of invariant polynomials. In other words $\Bbb{C}\left[\mathfrak{g}\right]^G=\Bbb{C}[\tau_2,..,\tau_n].$
\end{cor*}
\begin{proof}
If there a polynomial function $P:\Bbb{C}^{n-1} \to \Bbb{C}$ such that $P(\sigma_2,...,\sigma_n)=0$, then for all $(z_2,...,z_{n}) \in \Bbb{C}^{n-1}$ , consider the polynomial equation
$$
T^n-0T^{n-1}+z_2T^{n-2}-z_3T^{n-2}+...+\left( -1 \right) ^nz_n=0,
$$
the equation have $n$ solution $t_1,...,t_n $, by Vieta theorem so the equations 
$$
\sigma _1\left( x_1,...,x_n \right) =0,\sigma _2\left( x_1,...,x_n \right) =z_2,...,\sigma _n\left( x_1,...,x_n \right) =z_n
$$
have a solution $\left( x_1,...,x_n \right) =\left( t_1,...,t_n \right) $, this means
$$
P\left( z_2,...,z_n \right) =0,\forall \left( z_2,...,z_n \right) \in \mathbb{C} ^{n-1}.
$$
Thus $P=0$ . We have proved that the generators $\sigma_2,...,\sigma_n$ are algebraically independent. Similarly, we can prove $\tau_2,..,\tau_n$ are algebraically independent using the characteristic polynomial.
\end{proof} 

\begin{prop}[Type $B_n$ Case]
 When $G=\mathrm{SO}_{2n+1}(\Bbb{C})$, the induced map is surjective. In addition, 
$$
\mathbb{C} \left[ \mathfrak{h} \right] ^W=\mathbb{C} \left[ \sigma _1,...,\sigma _n \right] ,\quad \sigma _k\left( x_1,...,x_n \right) :=\sum_{i_1<...<i_k}{\left( x_{i_1}...x_{i_k} \right) ^2}
$$
\end{prop}
\begin{proof}
We have known that in this case $\mathfrak{g=so}_{2n + 1}(\mathbb{C})$ and one of Cartan subalgebra is
$$\mathfrak{h} = \mathrm{span}\{E_{j} - E_{n+j} : j = 1, \ldots, n\} = \{\mathrm{diag}(h, -h, 0) : h \in \mathbb{C}^n\}.$$
And the Weyl group $W\simeq S_n \ltimes \mathbb{Z}^{n}$ with natural action ($e_j:=E_{jj}-E_{n+j,n+j}$) generated by 
$$\sigma _{\varepsilon_k-\varepsilon_{k+1}}\left( e_{j} \right) =e_{\pi \left( j \right)},\pi =\left( k,k+1 \right) \in S_n;$$
$$
\sigma _{\varepsilon _n}\left( e_n\right) =-e_n;
$$
Let us identify $x\in  \mathfrak{h}$ as $x=(x_1,...,x_n)\in \Bbb{C}$, then $$
\sigma _{\varepsilon _n}P\left( x_1,...,x_n \right) =P\left( x_1,...,-x_n \right) ;
$$
$$
\sigma _{\varepsilon _k-\varepsilon _{k+1}}P\left( x_1,...,x_n \right) =P\left( x_{\pi \left( 1 \right)},...,x_{\pi \left( n \right)} \right) ,\pi =\left( k,k+1 \right) \in S_n.
$$
So $$
\mathbb{C} \left[ \mathfrak{h} \right] ^W\subseteq \mathbb{C} \left[ x_{1}^{2},...,x_{n}^{2} \right] ^{S_n}=\mathbb{C} \left[ \sigma _1,...,\sigma _n \right] ,\quad \sigma _k\left( x_1,...,x_n \right) :=\sum_{i_1<...<i_k}{\left( x_{i_1}...x_{i_k} \right) ^2}.
$$
In addition, the generators $\sigma_i$ are coefficients of (with $h=$diag$(x_1,...,x_n)$ )
$$
\det \left( tI-\left( \begin{matrix}
	h&		&		\\
	&		-h&		\\
	&		&		0\\
\end{matrix} \right) \right) =t\left( t^2-x_{1}^{2} \right) ...\left( t^2-x_{n}^{2} \right) .
$$
Similarly 
$$
\det \left( tI-\left( \begin{matrix}
	a&		b&		y\\
	c&		-a^T&		z\\
	y&		z&		f\\
\end{matrix} \right) \right) =t\sum_{k=0}^{n}{\left( -1 \right) ^{k}\tau _k}t^{2n-2k}
$$
is $G$-invariant and we have $\tau_k$ is $G$-invariant such that $\tau_{k}|\mathfrak{h}=\sigma_{k}$. Thus, the induced map is surjective. 
\end{proof}


\begin{prop}[Type $C_n$ case]
When $G=\mathrm{sp}_{2n}(\Bbb{C})$, the induced map is surjective. In addition,
$$
\mathbb{C}\left[ \mathfrak{h} \right] ^W=\mathbb{C} \left[ \sigma _1,...,\sigma _n \right] ,\quad \sigma _k\left( x_1,...,x_n \right) :=\sum_{i_1<...<i_k}{\left( x_{i_1}...x_{i_k} \right) ^2}.
$$
\end{prop}
\begin{proof}
  It is entirely similar to the proof of type $B_n$ case.  
\end{proof}


\begin{prop}[Type $D_n$ case]
 When $G=\mathrm{SO}_{2n}(\Bbb{C})$, the induced map is surjective. In addition,
$$
\mathbb{C} \left[ \mathfrak{h} \right] ^W=\mathbb{C} \left[ \sigma _1,...,\sigma _{n-1},x_1x_2..x_n \right] ,\quad \sigma _k\left( x_1,...,x_n \right) :=\sum_{i_1<...<i_k}{\left( x_{i_1}...x_{i_k} \right) ^2}.
$$
\end{prop}
\begin{proof}
 Let us identify $x\in  \mathfrak{h}$ as $x=(x_1,...,x_n)\in \Bbb{C}$, then 
$$
\sigma_{\varepsilon_n+\varepsilon_{n-1}}P\left( x_1,...,x_{n-1},x_n \right) =P\left( x_1,...,-x_{n-1},-x_n \right) ;
$$
$$
\sigma _{\varepsilon _k-\varepsilon _{k+1}}P\left( x_1,...,x_n \right) =P\left( x_{\pi \left( 1 \right)},...,x_{\pi \left( n \right)} \right) ,\pi =\left( k,k+1 \right) \in S_n.
$$
So $$
\mathbb{C} \left[ x_{1}^{2},...,x_{n}^{2} \right] ^{S_n}=\mathbb{C} \left[ \sigma _1,...,\sigma _n \right]\subseteq  \mathbb{C} \left[ \mathfrak{h} \right] ^W ,\quad \sigma _k\left( x_1,...,x_n \right) :=\sum_{i_1<...<i_k}{\left( x_{i_1}...x_{i_k} \right) ^2}.
$$
While $x_1x_2...x_n$ are also a invariant polynomial, we can vertify that $\mathbb{C} \left[ \mathfrak{h} \right] ^W=\mathbb{C} \left[ \sigma _1,...,\sigma _{n-1},x_1x_2..x_n \right] $. 
In addition, the generators $\sigma_i\,(i\le n-1)$ are coefficients of (with $h=$diag$(x_1,...,x_n)$ )
$$
\det \left( tI-\left( \begin{matrix}
	h&		\\
	&		-h\\
\end{matrix} \right) \right) =\left( t^2-x_{1}^{2} \right) ...\left( t^2-x_{n}^{2} \right) .
$$
Similarly, using the coefficient of characteristic polynomial, we have $\tau_i(1\le n-1)$ are $G$- invariant polynomials such that $\tau_i|_{\mathfrak{h}}=\sigma_i (1\le n-1)$.  According to \ref{lem:1}  , we know $\exists u,\forall x\in \mathfrak{g},(uxu^{-1})^{\top} =-uxu^{-1}$, and the determinant of skew-sysmmetric matrix is the square of Pfaffian. So $\det x=\mathrm{Pf}\left( uxu^{-1} \right) ^2$. Thus $$
\mathrm{Pf}\left( uxu^{-1} \right) =\sqrt{\det x}\in \mathbb{C} \left[ \mathfrak{g} \right] ^G,
$$
and its restriction on $\mathfrak{g}$ is $(\sqrt{-1})^nx_1x_2...x_n $. 
\end{proof}

\begin{defn}[rank and the regular element]
Let us denote the space $$\mathfrak{g} _{x}^{a}:=\mathfrak{g} ^a\left( \mathrm{ad}\,x \right) :=\bigcup_{n\geqslant 1}{\mathrm{Ker}\left( a\mathrm{id}-\mathrm{ad}\,x \right) ^n}.$$
    \begin{enumerate}
    \item The number
$$\mathrm{rank}(\mathfrak{g}) := \min \left\{ \dim \mathfrak{g}^{0}(\mathrm{ad} \, x) \mid x \in \mathfrak{g} \right\}.$$
    is called the \textbf{rank} of $\mathfrak{g}$. 
   \item An element $x \in \mathfrak{g}$ is called \textbf{regular} if $\dim \mathfrak{g}^{0}(\mathrm{ad} \, x)=\mathrm{rank}(\mathfrak{g})$.
    \item We write $\mathfrak{g}^r$ for the set of regular elements in $\mathfrak{g}$ and $\mathfrak{g}^s$ for the set of semisimple elements in $\mathfrak{g}$.
\end{enumerate}
\end{defn}


\begin{rem}
Since $\dim \mathfrak{g}^{0}(\mathrm{ad}\, x)$ is the multiplicity of $0$ as a root of the characteristic polynomial
$$\det(\mathrm{ad} \, x - t \mathrm{id}_\mathfrak{g})=\sum_{k = 0}^{n} p_k(x)t^k .$$
    of $\mathrm{ad}\,x$, we have
    $$\mathrm{rank}(\mathfrak{g})=\min \left\{ k \in \mathbb{N} \mid p_k \neq 0 \right\}.$$  
Let $\{e_j,1\le j\le n\}$ be a basis of $\mathfrak{g}$ and $[e_i,e_j]=\sum_{k=1}^n a_{i,j}^ke_k,\quad
x=\sum_{i=1}^n{x_ie_i}$ then
$$
\mathrm{ad}\,x\left( e_j \right) =\left[ x,e_j \right] =\sum_{i=1}^n{x_i[e_i,e_j]}=\sum_{i=1}^n{\sum_{k=1}^n{a_{i,j}^{k}}x_ie_k}=\sum_{k=1}^n{\left( \sum_{i=1}^n{a_{i,j}^{k}}x_i \right) e_k}.
$$
So the entries of the matrix of $\mathrm{ad}\,x$ are linear combination of $x_i$, which implies that all the functions $p_k: \mathfrak{g} \to \Bbb{C}$  are polynomials.

Since the characteristic polynomial is invariant under adjoint action, we can immediately conclude the following.
\end{rem}

\begin{lem}[Chap.6 in \cite{SGLG} ]\label{lem:3}
The polynomial $p_r(x)\in \Bbb{C}[\mathfrak{g}]^G$ and 
$$
\mathfrak{g} ^r=\left\{ x\in \mathfrak{g} |p_r\left( x \right) \ne 0 \right\} .
$$
In addtion, $\mathfrak{g} ^r$ is an open, $G$-stable,  dense and connected subset of $\mathfrak{g}$. 
\end{lem}
\begin{proof}
The proof of this lemma about dense and connected in the reference book  \cite{SGLG} uses the results of Exercise 6.1.2 and 6.1.2. Next, I will prove the results of these two exercises.

To show that the set is dense in $\mathfrak{g}$, equivalently, $Z:=\{x\in \Bbb{C}^n:p(x)=0\}$ has no interior points, where $p$ is a non-constant polynomial. Otherwise, if $p(x)=0, \forall x \in B(x_0,r)$, then using the Taylor expansion at $x=x_0$, $$
p\left( x \right) =\sum_{|\alpha |=k}^{\mathrm{deg}\,p}{\frac{\partial ^{\alpha}p\left( x_0 \right)}{\alpha !}\left( x-x_0 \right) ^{\alpha}},\quad \alpha =\left( \alpha _1,..,\alpha _n \right) ,
$$
we can obtain $p=0$. 

To show that the set $X=\mathfrak{g} ^r = \{ x \in \mathfrak{g} \mid p(x) \ne 0 \}$ is connected, we consider the properties of complex algebraic varieties and their complements.  

First, recall that $\mathbb{C}^n$ is homeomorphic to $\mathbb{R}^{2n}$, which is connected. The zero set of a non-constant polynomial $p(x)$, denoted $Z = \{ x \in \mathbb{C}^n \mid p_r(x) = 0 \}$, is a proper algebraic subset of $\mathbb{C}^n$. Specifically, $Z$ is a complex hypersurface of complex dimension $n - 1$, which corresponds to real dimension $2n - 2$ in $\mathbb{R}^{2n}$. 

In topology, removing a subset of codimension greater than 1 from a manifold does not disconnect the manifold. Since $Z$ has real codimension 2 in $\mathbb{R}^{2n}$, removing $Z$ from $\mathbb{R}^{2n}$ does not disconnect it. Therefore, the complement $X = \mathbb{C}^n \setminus Z$ remains connected.

Therefore, the set $X = \mathfrak{g} ^r $ is connected.
\end{proof}

\begin{lem}[Chap.6 in \cite{SGLG} ]
 Let $\mathfrak{g}$ be semisimple Lie algebra over $\Bbb{C}$ and $\mathfrak{h}$ be a Cartan subalgebra , then 
\begin{enumerate}
    \item If $
\Delta =\Delta \left( \mathfrak{g} ,\mathfrak{h} \right) 
$ be a base, then
$$
\mathfrak{g} ^r\cap \mathfrak{h} =\mathfrak{h} \backslash \bigcup_{\alpha \in \Delta}{\mathrm{Ker}}\alpha \ne\emptyset.
$$
\item  For all $x\in \mathfrak{g}^r$, the space $\mathfrak{g}^0_{x}$  form a Cartan subalgebra.
\item For all $x\in\mathfrak{g}^r\cap \mathfrak{h}$ , we have $\mathfrak{h}=\mathfrak{g}^0_{x}$. 
\item Every regular element in $\mathfrak{g}$ is simisimple. In other words
$$
\mathfrak{g} ^r\cap \mathfrak{g} ^s=\mathfrak{g} ^r.
$$
\end{enumerate}
\end{lem}
\begin{cor*}
    Let $H$ be the set of  Cartan subalgebras of  $\mathfrak{g}$, then $$
\mathfrak{g} ^r\subseteq \bigcup_{\mathfrak{h} \in H}{\mathfrak{h}}.
$$\end{cor*}
\begin{proof}
If $x\in \mathfrak{g}^r$, then $x\in \mathfrak{g}^0_x=\mathfrak{h}$, which is a Cartan subalgebra.
\end{proof}

\begin{lem}[Chap.6 in \cite{SGLG} ]
$\mathrm{Inn}(\mathfrak{g})$ acts transitively on the set of Cartan subalgebra of $\mathfrak{g}$.   In other words, if we fixed a Cartan subalgebra, then 
$$
H=\left\{ g\mathfrak{h} |g\in \mathrm{Inn}\left( \mathfrak{g} \right) \right\} \Rightarrow \mathfrak{g} ^r\subseteq \bigcup_{g\in \mathrm{Inn}\left( \mathfrak{g} \right)}{g\mathfrak{h}}\subseteq \bigcup_{g\in G}{g\mathfrak{h}g^{-1}}.
$$
\end{lem}


\begin{prop}[Injective]
 For general complex connected Lie group, the induced map  is an injection.  
\end{prop}
\begin{proof}
 Let $P\in \Bbb{C}[\mathfrak{h}]^W$ such that $P|_{\mathfrak{h}}=0$, then we have
$$
\forall x\in \bigcup_{g\in G}{g\mathfrak{h}g^{-1}}, \,P(x)=0.
$$
By the lemma above, we have 
$$
\forall x\in \mathfrak{g}^r,P(x)=0, \quad \mathfrak{g}^r
=\{p_r(x)\ne 0|x\in \mathfrak{g}\}.
$$
But \ref{lem:3} we have known that $\mathfrak{g}^r$ dense in $\mathfrak{g}$,  which implies $P=0$. 

In summary , we have proved the kernel of the induced map is 0. 
\end{proof}

\vspace{0.5cm}
\section{Conclusion and Future Work}
I have proved that the Chevalley restriction theorem holds true for the case of classical simple Lie groups, and worked out the concrete generators of the invariant polynomial functions in $\Bbb{C}[\mathfrak{g}]^G$ and $\Bbb{C}[\mathfrak{h}]^W$.

In the future, I will strive to calculate the generators in the Chevalley restriction theorem for exceptional complex simple Lie algebras, and prove that  the induced map is also a surjection for general complex semisimple Lie groups.



\begin{thebibliography}{99}
\bibitem{GTM09}
James E.Humphreys,* Introduction to Lie Algebra and Representation Theory GTM 9*,Springer, 1972. Chap 1-11.
\bibitem{SGLG}
Joachim Hilgert and Karl-Hermann Neeb:Structure and Geometry
 of Lie Groups\\
https://softbank.iust.ac.ir/MathBooks/N/Neeb%20and%20Hilgert%20-%20Structure%20and%20Geometry%20of%20Lie%20Groups.pdf
\bibitem{RTCG}
Neil Chriss and Victor Ginzburg, *Representation Theory and Complex Geometry*, Springer, 1997. Section 3.1.37.\\
https://link.springer.com/book/10.1007/978-0-8176-4938-8
\bibitem{LA}
Zhexian Wan,*Lie Algebra(second ediction)*,Higher Education Press,2013,Section 6.2-6.3
\bibitem{CN}
Chen, T. H. and Ng\^o, B. C.: On the Hitchin fibration for algebraic surfaces. arXiv: 1711.02592v2 (2018).
\end{thebibliography}

{\color{red}try to use bibtex for references}

\bigskip
\noindent\small{\textsc{Department of Mathematics, South China University of Technology}\\
\emph{E-mail address}:  \texttt{lmnjkljh@sina.com}

\end{document}

